% P15.34: Feller II VII.1--VII.4, retaining its finite-difference and ramp route.
\chapter{Bernstein Approximation, Hausdorff Moments and Exchangeability}
\label{sec:p15-feller-bernstein-hausdorff}

A probability distribution on a bounded interval is determined by its
moments. Feller proves this by turning finite differences into probability
weights, constructing a distribution function with continuous ramps, and
then applying the result to exchangeable observations. This chapter retains
that route and its approximation premises
\cite[vol.~II, VII.1--VII.4, pp.~219--230]{Feller1971ProbabilityII}.
The moment proof below does not assume a measure representing its limiting
functional before that measure has been constructed.

\section{Concentration and the source approximation kernels}

\begin{hbtheorem}{thm:p15-feller-concentration}{proved=annex:p15.feller-concentration}{Feller's parameter-dependent concentration lemma}
For positive integers \(n\) and every parameter \(\theta\) in an interval \(I\), let
\(\nu_{n,\theta}\) be a real probability law with mean \(\theta\)
and finite variance \(v_n(\theta)\). If \(v_n(\theta)\to0\), then
for every bounded continuous \(u:\mathbb R\to\mathbb R\),
\[
                        \int u\,d\nu_{n,\theta}\longrightarrow u(\theta).
\]
Convergence is uniform on each compact parameter interval on which
\(v_n\to0\) uniformly. If \(u\) is uniformly continuous on all of
\(\mathbb R\), the same assertion holds on any parameter set with
uniformly vanishing variance. These are approximation conclusions,
not assertions that every finite-\(n\) expectation is already \(u(\theta)\).
\end{hbtheorem}

\begin{hbtheorem}{thm:p15-feller-kernel-examples}{proved=annex:p15.feller-kernel-examples}{The binomial, Poisson, gamma and plug-in examples}
For continuous \(u\) on \([0,1]\),
\[
 \sum_{j=0}^n u(j/n)\binom nj\theta^j(1-\theta)^{n-j}
                         \longrightarrow u(\theta)
\]
uniformly for \(0\leq\theta\leq1\).
Here \(n\) is a positive integer, and exponent-zero powers are one.

For bounded continuous \(u\) on \([0,\infty)\) and positive real
\(n\to\infty\),
\[
 e^{-n\theta}\sum_{j=0}^\infty u(j/n)\frac{(n\theta)^j}{j!}
                         \longrightarrow u(\theta)
\]
uniformly on compact subsets of \(0\leq\theta<\infty\).
For \(\theta>0\), the gamma approximation is
\[
 \frac1{\Gamma(n)}\int_0^\infty
       u(x)(nx/\theta)^{n-1}e^{-nx/\theta}\frac n\theta\,dx
                         \longrightarrow u(\theta).
\]
Its law has shape \(n\), rate \(n/\theta\), mean \(\theta\)
and variance \(\theta^2/n\). At \(\theta=0\), define the law as
\(\delta_0\); convergence is then uniform on compact nonnegative
parameter intervals as well. Integer \(n\) gives \(\Gamma(n)=(n-1)!\).

If independent observations have a common law with mean \(\theta\) and
finite variance, their sample mean \(M_n\) gives
\(E u(M_n)\to u(\theta)\) for bounded continuous \(u\).
Thus the corresponding plug-in estimator is asymptotically unbiased
under these hypotheses. These are the four source examples in VII.1.
\end{hbtheorem}

\section{Finite differences and Bernstein polynomials}

\begin{hbdefinition}{def:p15-feller-differences}{Sequence differences and normalized function differences}
For a real sequence \((a_i)\), put
\(\Delta a_i=a_{i+1}-a_i\) and
\(\Delta^0a_i=a_i\), with higher powers by iteration.
A finite sequence permits only indices for which all terms exist.
For a real function \(u\) and \(h>0\), the source's normalized
difference is
\[
                 \Delta_h u(x)=\frac{u(x+h)-u(x)}h,
                 \qquad \Delta_h^0u(x)=u(x).
\]
Higher powers keep the same \(h\) and require the full grid
\(x,x+h,\ldots,x+rh\) in the domain.
The unadorned \(\Delta\) acts on sequences; \(\Delta_h\) is the
normalized function operator.
\index{finite differences}
\end{hbdefinition}

\begin{hbtheorem}{thm:p15-feller-differences}{proved=annex:p15.feller-differences}{The complete finite-difference reciprocity}
For integers \(r\geq0\),
\[
 \Delta^r a_i=\sum_{j=0}^r\binom rj(-1)^{r-j}a_{i+j},
 \qquad
 \Delta_h^r u(x)=h^{-r}\sum_{j=0}^r\binom rj(-1)^{r-j}u(x+jh).
\]
For two real sequences \(a,c\), and integers \(i,N\geq0\),
\[
 \sum_{r=0}^N\binom Nr c_r\Delta^r a_i
   =\sum_{j=0}^N\binom Nj a_{i+j}(-1)^{N-j}\Delta^{N-j}c_j.
\]
In particular
\[
                c_k=\sum_{j=0}^N\binom Nj(-1)^{N-j}
                                            \Delta^{N-j}c_{k+j}.
\]
The latter inversion holds for every \(k,N\geq0\).
Taking \(c_r=\theta^r\) gives the source's binomial reciprocity,
including \(\theta=0,1\) with the stated power convention.
\end{hbtheorem}

\begin{hbdefinition}{def:p15-feller-bernstein}{The explicit approximation polynomial}
For \(u:[0,1]\to\mathbb R\) and integer \(n\geq1\), define
\[
 B_{n,u}(\theta)=
      \sum_{j=0}^n u(j/n)\binom nj\theta^j(1-\theta)^{n-j}.
\]
For continuous \(u\), let
\(\|u\|_\infty=\max_{[0,1]}|u|\), and
\(\omega_u(\delta)=\sup\{|u(x)-u(y)|:x,y\in[0,1],\,|x-y|\leq\delta\}\).
\index{Bernstein polynomial}
\end{hbdefinition}

\begin{hbtheorem}{thm:p15-feller-bernstein}{proved=annex:p15.feller-bernstein}{The two representations and uniform approximation}
Writing \(h=1/n\), the same polynomial is
\[
                   B_{n,u}(\theta)
                       =\sum_{r=0}^n\binom nr(h\theta)^r\Delta_h^ru(0).
\]
If \(u\) is continuous on \([0,1]\), then \(B_{n,u}\to u\)
uniformly, with the explicit bound
\[
 \|B_{n,u}-u\|_\infty
        \leq\omega_u(\delta)+\frac{\|u\|_\infty}{2n\delta^2}
                                         \quad(\delta>0).
\]
If \(u\) is \(L\)-Lipschitz, the bound \(L/(2\sqrt n)\) also holds.
Affine change of variable supplies uniform polynomial approximation on
every finite interval \([a,b]\), \(a<b\).
\end{hbtheorem}

\begin{hblemma}{thm:p15-feller-positive-series}{proved=annex:p15.feller-positive-series}{The positive-coefficient limit premise}
Let
\(P_n(x)=\sum_{r\geq0}a_{n,r}x^r\), where \(a_{n,r}\geq0\) and
\(\sum_r a_{n,r}\leq M<\infty\).
If \(P_n(x)\to u(x)\) for every \(0\leq x<1\), then there are
unique \(a_r\geq0\), with \(\sum_r a_r\leq M\), such that
\[
                      a_{n,r}\to a_r,\qquad
                      u(x)=\sum_{r=0}^\infty a_rx^r\quad(0\leq x<1).
\]
If \(u\) extends continuously to \(1\), then \(\sum_r a_r=u(1)\).
In particular a limit of probability generating functions that is
continuous at \(1\) with value one is a probability generating function.
The continuity requirement excludes loss of mass, as \(P_n(x)=x^n\)
shows.
\end{hblemma}

\begin{hbtheorem}{thm:p15-feller-absolute-monotone}{proved=annex:p15.feller-absolute-monotone}{The exact three-way absolute-monotonicity criterion}
For a continuous \(u:[0,1)\to\mathbb R\), the following are equivalent:
\begin{enumerate}
\item \(u(x)=\sum_{r=0}^\infty p_rx^r\) on \([0,1)\), with
      \(p_r\geq0\) and the series finite at every \(x<1\).
\item \(u\) is infinitely differentiable on \((0,1)\) and
      \(u^{(r)}(x)\geq0\) for every integer \(r\geq0\) and \(0<x<1\).
      Order zero includes nonnegativity of \(u\).
\item For every integer \(n\geq1\),
      \(\Delta_{1/n}^{\,k}u(0)\geq0\) for \(0\leq k\leq n-1\).
\end{enumerate}
The grid condition never evaluates \(u(1)\).
The proof retains the Bernstein positive-coefficient construction,
using subinterval grids to justify the stated \(k<n\) condition
without adding an extra endpoint hypothesis.
The coefficients are unique. The function is a probability generating
function exactly when \(u(1-)=1\); this requires the boundary limit,
not a separately assigned value at \(1\).
The zero function and unbounded absolutely monotone functions are included.
\index{absolute monotonicity}
\end{hbtheorem}

\section{From complete monotonicity to a probability measure}

\begin{hbdefinition}{def:p15-feller-complete-monotone}{Moment sequences and their finite probability weights}
A real sequence \((c_k)_{k\geq0}\) is completely monotone when
\[
                         (-1)^r\Delta^r c_k\geq0
                                            \quad(k,r\geq0).
\]
For a probability \(\mu\) on \([0,1]\), its moments are
\(c_k=\int_{[0,1]}x^k\,\mu(dx)\), with \(x^0=1\), including \(x=0\).
For a completely monotone sequence with \(c_0=1\), and integers
\(n\geq1\), \(0\leq j\leq n\), put
\[
 w_{n,j}=\binom nj(-1)^{n-j}\Delta^{n-j}c_j,\qquad
 \mu_n=\sum_{j=0}^n w_{n,j}\delta_{j/n}.
\]
The following result proves that these are probability measures,
without presupposing an original measure for the sequence.
\index{complete monotonicity!sequences}
\end{hbdefinition}

\begin{hbtheorem}{thm:p15-feller-hausdorff-weights}{proved=annex:p15.feller-hausdorff-weights}{The source's constructive weights and polynomial moments}
A probability on \([0,1]\) has \(c_0=1\) and
\[
                      (-1)^r\Delta^r c_k
                         =\int x^k(1-x)^r\,\mu(dx)\geq0.
\]
Conversely, for a completely monotone sequence with \(c_0=1\),
the weights \(w_{n,j}\) are nonnegative and sum to one.
Write \(E_nu=\int u\,d\mu_n\).
For \(h=1/n\) and any grid-defined \(u\),
\[
                 E_nu=\sum_{r=0}^n\binom nr h^r c_r\Delta_h^ru(0).
\]
For every polynomial \(u\) of degree \(d\),
\[
                 E_nu\longrightarrow
                   \sum_{r=0}^d c_r\,\frac{u^{(r)}(0)}{r!}.
\]
In particular the \(k\)th moments of \(\mu_n\) tend to \(c_k\).
\end{hbtheorem}

\begin{hbtheorem}{thm:p15-feller-ramp-measure}{proved=annex:p15.feller-ramp-measure}{The continuous-ramp construction, including endpoint masses}
For the preceding weights, \(E_nu\) converges for every continuous
\(u\) on \([0,1]\). Denote the limit by \(L(u)\).
It is positive and linear, with \(L(1)=1\) and
\(|L(u)|\leq\|u\|_\infty\).

For real \(t\), \(\epsilon>0\), and \(0\leq x\leq1\), put
\[
 u_{t,\epsilon}(x)=
 \begin{cases}
 1,&x\leq t,\\
 (t+\epsilon-x)/\epsilon,&t<x<t+\epsilon,\\
 0,&x\geq t+\epsilon.
 \end{cases}
 \qquad
 U_\epsilon(t)=L(u_{t,\epsilon}),\quad
 F(t)=\inf_{\epsilon>0}U_\epsilon(t).
\]
Then \(F\) is a right-continuous nondecreasing CDF, zero for \(t<0\)
and one for \(t\geq1\). In particular \(F(0)\) need not be zero.
It defines a probability \(\mu\) supported on \([0,1]\).
The CDFs \(F_n\) of \(\mu_n\) converge to \(F\) at every
continuity point of \(F\). The proof derives right continuity directly
from the ramps before invoking the CDF measure construction.
\end{hbtheorem}

\begin{hbtheorem}{thm:p15-feller-hausdorff}{proved=annex:p15.feller-hausdorff}{Hausdorff's theorem and the full source inversion}
A sequence is the moment sequence of a unique probability on \([0,1]\)
exactly when it is completely monotone and \(c_0=1\).
For that probability, at every continuity point \(t\) of its CDF,
\[
                       \sum_{\substack{0\leq j\leq n\\j\leq nt}}
                            \binom nj(-1)^{n-j}\Delta^{n-j}c_j
                                  \longrightarrow F(t).
\]
The construction also satisfies \(L(u)=\int u\,d\mu\) for every
continuous \(u\), and for every integer \(r\geq1\),
\[
                \int x^r\,\mu(dx)
                    =1-r\int_0^1 t^{r-1}F(t)\,dt.
\]
No atom is discarded at either endpoint.
In particular
\[
 \mu(\{1\})=\lim_{n\to\infty}c_n,\qquad
 \mu(\{0\})=\lim_{n\to\infty}(-1)^n\Delta^nc_0.
\]
A completely monotone sequence with finite \(c_0\geq0\) similarly
represents a unique finite measure of total mass \(c_0\);
\(c_0=0\) forces the zero sequence and zero measure.
\index{Hausdorff moment theorem}\index{moment problem!bounded interval}
\end{hbtheorem}

\begin{hbtheorem}{thm:p15-feller-moment-calculations}{proved=annex:p15.feller-moment-calculations}{Explicit inversion, affine intervals and a finite-prefix obstruction}
\begin{enumerate}
\item \(c_k=1/(k+1)\) gives the uniform law, with
      \(w_{n,j}=1/(n+1)\).
\item \(c_0=1,c_k=a\) for \(k\geq1\), \(0\leq a\leq1\),
      gives \((1-a)\delta_0+a\delta_1\); the finite weights already
      give this exact law at every \(n\geq1\).
\item \(c_k=t_0^k\), \(0<t_0<1\), gives \(\delta_{t_0}\).
      Its weights are binomial probabilities. At the discontinuity
      \(t_0\), \(F_n(t_0)\to1/2\), whereas \(F(t_0)=1\).
      The continuity-point qualification is essential.
\item For finite \(a<b\), define
      \(d_k=(b-a)^{-k}\sum_{j=0}^k\binom kj(-a)^{k-j}c_j\).
      A sequence with \(c_0=1\) represents a probability on \([a,b]\)
      exactly when \(d\) is completely monotone.
      For \(a=b\), the criterion is \(c_k=a^k\), with the order-zero
      value one.
\item The prefix \(c_0=1,c_1=1/2,c_2=1/10\) satisfies every
      available complete-monotonicity inequality with \(k+r\leq2\),
      but cannot be the moments of a probability.
      Finite-prefix difference tests do not establish infinite extendibility.
\end{enumerate}
\end{hbtheorem}

\section{Unbounded moment problems and the source exchangeability proof}

\begin{hbtheorem}{thm:p15-feller-lognormal-moments}{proved=annex:p15.feller-lognormal-moments}{The source's lognormal counterexample and sufficient moment bounds}
For \(x>0\), put
\[
 f(x)=\frac{1}{\sqrt{2\pi}x}e^{-(\log x)^2/2},\qquad
 f_a(x)=f(x)\{1+a\sin(2\pi\log x)\},\quad -1\leq a\leq1,
\]
and set both densities to zero for \(x\leq0\).
They are distinct probability densities as \(a\) varies, and all
their integer moments equal \(e^{k^2/2}\), \(k\geq0\).
The annex gives Feller's shifted odd-integrand proof.

For a real probability law with moments \(m_k\), the source's sufficient
condition
\[
                 \sum_{n=0}^\infty\frac{m_{2n}s^n}{(2n)!}<\infty
                               \quad\text{for some }s>0
\]
implies moment determinacy. The stronger Carleman condition
\(\sum_{n\geq1}m_{2n}^{-1/(2n)}=\infty\) has its complete proof in
\cref{thm:p15-criteria-carleman}. A zero positive-order even moment
means the point mass at zero and is treated separately.
\end{hbtheorem}

\begin{hbtheorem}{thm:p15-feller-binary-mixture}{proved=annex:p15.feller-binary-mixture}{The complete finite-difference de Finetti argument}
Let \(X_1,X_2,\ldots\) be an infinite exchangeable binary sequence.
Put \(c_0=1\),
\(c_n=P(X_1=\cdots=X_n=1)\), and
\[
 p_{k,n}=P(X_1=\cdots=X_k=1,\ X_{k+1}=\cdots=X_n=0)
                      \quad(0\leq k\leq n).
\]
Then
\[
                         p_{k,n}=(-1)^{n-k}\Delta^{n-k}c_k.
\]
Thus \(c\) is completely monotone and determines a unique probability
\(\lambda\) on \([0,1]\), giving
\[
 p_{k,n}=\int\theta^k(1-\theta)^{n-k}\,\lambda(d\theta),\qquad
 P\!\left(\sum_{j=1}^nX_j=k\right)=\binom nk p_{k,n}.
\]
Conversely any such \(\lambda\) constructs an exchangeable sequence
with these probabilities. This is Feller's moment route, independent
of the reverse conditional-product proof.
\index{exchangeability!moment representation}
\end{hbtheorem}

\begin{hbtheorem}{thm:p15-feller-polya-moments}{proved=annex:p15.feller-polya-moments}{The source's urn identification by its all-success moments}
For the urn with positive black weight \(b\), red weight \(r\)
and same-color reinforcement \(c>0\), let \(X_j\) indicate black.
Its all-success probabilities are
\[
 c_n=\frac{b(b+c)\cdots(b+(n-1)c)}
             {(b+r)(b+r+c)\cdots(b+r+(n-1)c)}
     =\frac{\Gamma(b/c+n)\Gamma((b+r)/c)}
             {\Gamma((b+r)/c+n)\Gamma(b/c)}.
\]
The empty product is one. These are precisely the moments of
\(\operatorname{beta}(b/c,r/c)\), so Hausdorff uniqueness identifies
the mixing law. The beta integral then supplies every finite word
and count probability, agreeing with the earlier direct urn calculation.
\end{hbtheorem}

\begin{hbtheorem}{thm:p15-feller-finite-alphabet}{proved=annex:p15.feller-finite-alphabet}{The stated finite-alphabet generalization}
For an infinite exchangeable sequence taking values in
\(\{1,\ldots,d\}\), \(d\geq1\), there is a unique probability \(\lambda\)
on
\(\Delta_d=\{p\in[0,1]^d:\sum_i p_i=1\}\) such that a specified word
with counts \(n_1,\ldots,n_d\), \(\sum_i n_i=n\), has probability
\[
                              \int_{\Delta_d}\prod_{i=1}^dp_i^{n_i}\,
                                                       \lambda(dp).
\]
The count event has the additional multinomial factor
\(n!/\prod_i n_i!\). The empirical frequency vector converges almost
surely to a vector with law \(\lambda\).
The annex gives the explicit bridge from the already proved
standard-Borel directing-law theorem to this source generalization.
\end{hbtheorem}

\begin{hbtheorem}{thm:p15-feller-finite-exchangeability}{proved=annex:p15.feller-finite-exchangeability}{The two source finite-sequence obstructions}
Choose uniformly from the six weak compositions of two balls into
three labeled cells, and let \(X_i\) indicate whether cell \(i\) is
occupied. This vector is exchangeable, with
\(c_0=1,c_1=1/2,c_2=1/6,c_3=0\).
It has no infinite exchangeable extension and no identical-independent
mixture: any completely monotone continuation would have \(c_j=0\)
for \(j\geq3\), giving \(\Delta^4c_1=-1/6\).

For i.i.d. real \(Z_1,\ldots,Z_n\), \(n\geq3\), with a
nondegenerate common law, meaning a law not concentrated at one point, put
\(Y_k=Z_k-n^{-1}\sum_{i=1}^nZ_i\), \(1\leq k\leq n-1\).
The vector is exchangeable and has no identical-independent mixture.
This holds without a moment hypothesis. When in addition
\(0<\operatorname{Var}(Z_1)=\sigma^2<\infty\), its covariance
is \(\operatorname{Cov}(Y_i,Y_j)=-\sigma^2/n<0\) for \(i\ne j\).
For \(n=2\) the retained vector has only one coordinate and admits
a mixture of point masses; degenerate input likewise gives a degenerate
mixture. The dimension and nondegeneracy qualifications are essential.
The annex gives the full bounded-transform argument and the
additional finite-variance calculation.
\end{hbtheorem}

\begin{hbtheorem}{thm:p15-feller-approximation-boundaries}{proved=annex:p15.feller-approximation-boundaries}{Why the approximation and positivity hypotheses matter}
\begin{enumerate}
\item Mean-zero laws with masses \(1/(2n^3)\) at each of \(-n,n\)
      and the remaining mass at zero have variance \(1/n\to0\),
      but their expectation of \(u(x)=x^4\) is \(n\).
\item There is a bounded continuous \(u\) on \(\mathbb R\) for which
      the laws \(\tfrac12\delta_{\theta-1/n}+\tfrac12\delta_{\theta+1/n}\)
      have uniformly vanishing variance yet their \(u\)-expectations
      do not converge uniformly over all real \(\theta\).
      A complete explicit bump-function example is given in the annex.
\item The function \(u=-1\) has nonnegative derivatives of every
      positive order but fails the absolute-monotonicity criterion
      because order zero is negative.
      The nonnegative function \(u(x)=x-x^2\) on \([0,1]\)
      fails its second-difference condition.
\item \(u(x)=(1-x)^{-1}\) is absolutely monotone on \([0,1)\)
      with coefficients all one, but is not a probability generating
      function. The zero function satisfies the analytic criterion
      without a permissible division by its value at a positive point.
\end{enumerate}
\end{hbtheorem}

\begin{proofboundary}{def:p15-feller-moment-scope}{The completed Feller route and its separate continuations}
The supplied route is Feller Volume II, VII.1--VII.4: the concentration
lemma and four examples, finite-difference reciprocity, both Bernstein
theorems, the complete Hausdorff construction and inversion, the moment
counterexample, and the exchangeability theorem with all three source
examples. The residual example covers every nondegenerate real input law
for \(n\geq3\), and includes the dimension and degeneracy failure cases.

The Poisson semigroup and Laplace-inversion continuations in VII.5--VII.6,
the wider transform-continuity presentations, the separate original
Hausdorff and Heyde papers, and Feller's later analytic moment proof
retain their own source presentations. The finite-alphabet conclusion
above uses the explicit earlier directing-law bridge.
\end{proofboundary}
